% Drawing 121/03-2 -- folder 121, batch 03, archive page 41 (PDF page 42).
% Transcription: en haut à gauche, trois dessins l'un sous l'autre : un point « $V_0$ » ; une droite horizontale portant le point $V_0$ entre deux petits traits verticaux, « $V_1$ » de part et d'autre, deux petites flèches vers les traits et la légende « $\partial V_1 = V_{1,0}$ » ; une figure plane : la droite épaisse $V_1$ et son point $V_0$, bordée de deux traits parallèles qui contournent $V_0$ par deux arcs, les bandes marquées $V_{21}$, $V_{210}$, $V_{20}$, la région hachurée au-dessus et au-dessous marquée $V_2$.
% Kind: stratified neighbourhoods: a point, a line through it, a plane figure. Method 2 (schematic redraw from the facsimile). Status: partial.
% Notes: Three figures. (1) A point labelled V_0. (2) A horizontal line carrying the point V_0 between two short heavy vertical ticks, a small tick further left; V_1 written under each half of the line; two small arrows pointing up to the ticks, and under them the legend dV_1 = V_{1,0}. (3) The plane figure: a heavy middle line (V_1 on both halves) with the point V_0 and two small points on either side of it; a parallel line above and one below, each leaving its level to go round V_0 by an arc (the lower arc heavy and doubled), with a point at each end of each arc; labels V_{21}, V_{210}, V_{20}, V_{210}, V_{21} above and the same below; oblique hatching strokes above and below, and V_2 under the figure. Partial: the superscripts of V_0 (top) and of the two V_1 and of dV_1 are overwritten or struck scribbles and are left out; the hatching strokes are placed by eye (seven above, six below, as on the page); the arcs are drawn as circular arcs between their computed end points. The small arrow diagram at the right of the page (dV_1 = V_{10} -> V_0, and an arrow down to V_1) is a relation to be set as text and is not drawn. Redrawn 2026-11-12 (issue 128): to check again.
% Coordinates: 1 unit = 1 mm = 1/100 of the crop width (bbox in page fractions: [0.11, 0.09, 0.41, 0.31]).
% Redrawn by Claude (claude-fable-5-2), issue #128, 2026-11-12; not checked by a human; not compiled (no TeX in the session).
% Rights: reproduces part of the work -- behind RENDER_TRANSCRIPTS, never published by itself.
\documentclass[tikz,border=3pt]{standalone}
\usepackage{amsmath,amssymb}
\usetikzlibrary{patterns,arrows.meta}
\tikzset{pencil/.style={black!45}, struck/.style={black!35}, mpt/.style={circle,draw,inner sep=1.3pt}, pt/.style={circle,fill,inner sep=0.9pt}, opt/.style={circle,draw,fill=white,inner sep=1.1pt}}
\begin{document}
\begin{tikzpicture}[x=1mm,y=1mm,line cap=round,line join=round,font=\scriptsize]
% (1) the point V_0
\node[pt] at (51,-13.5) {};
\node[inner sep=1pt] at (56,-14) {\small $V_0$};
% (2) the line V_1 through V_0
\draw (19,-29.5) -- (78,-29.5);
\node[pt] at (51,-29.5) {};
\draw[thick] (45.5,-28) -- (45.5,-31);
\draw[thick] (56,-28) -- (56,-31);
\draw (27.5,-28.5) -- (27.5,-29.5);
\draw[-{Stealth[length=1.2mm]}] (47.5,-37) -- (46.5,-33.5);
\draw[-{Stealth[length=1.2mm]}] (52,-38) -- (55,-34);
\node[inner sep=1pt] at (29,-35) {\small $V_1$};
\node[inner sep=1pt] at (72,-35) {\small $V_1$};
\node[inner sep=1pt] at (53,-44) {\small $\partial V_1 = V_{1,0}$};
% (3) the plane figure
\coordinate (a1) at (51,-65.5);
\coordinate (a2) at (64,-65.5);
\coordinate (b1) at (51,-72);
\coordinate (b2) at (64,-72);
\draw (21,-66.5) -- (a1) to[out=80,in=180] (57.5,-60.5) to[out=0,in=100] (a2) -- (100,-65);
\draw (21,-73) -- (b1);
\draw[very thick] (b1) to[out=-80,in=180] (57.5,-76.5) to[out=0,in=-100] (b2);
\draw (b2) -- (100,-71.5);
\draw[very thick] (21,-70) -- (92,-69);
\node[opt] at (a1) {};
\node[opt] at (a2) {};
\node[pt,inner sep=1.3pt] at (b1) {};
\node[pt,inner sep=1.3pt] at (b2) {};
\node[pt,inner sep=0.6pt] at (51.5,-69.6) {};
\node[pt] at (56,-69.5) {};
\node[pt,inner sep=0.6pt] at (64,-69.4) {};
\foreach \p/\q in {(25,-66)/(34,-57), (37,-65.5)/(45.5,-56.5), (49,-61)/(58,-53.5), (62,-60)/(70.5,-52), (73,-64.5)/(82,-56.5), (85,-65)/(93,-58), (96,-65)/(100,-62)} {\draw[pencil] \p -- \q;}
\foreach \p/\q in {(27,-75)/(20,-88), (43,-74)/(34.5,-91), (55,-77)/(49,-84), (73,-76)/(66.5,-91), (88,-74)/(80.5,-89), (99,-74)/(90.5,-92)} {\draw[pencil] \p -- \q;}
\node[inner sep=1pt] at (35,-62.5) {\small $V_{21}$};
\node[inner sep=1pt] at (49,-62) {\small $V_{210}$};
\node[inner sep=1pt] at (59.5,-58) {\small $V_{20}$};
\node[inner sep=1pt] at (68,-62) {\small $V_{210}$};
\node[inner sep=1pt] at (89,-62.5) {\small $V_{21}$};
\node[inner sep=1pt] at (59,-67.3) {\small $V_0$};
\node[inner sep=1pt] at (32,-68) {\small $V_1$};
\node[inner sep=1pt] at (78,-67.3) {\small $V_1$};
\node[inner sep=1pt] at (34,-78) {\small $V_{21}$};
\node[inner sep=1pt] at (48.5,-78) {\small $V_{210}$};
\node[inner sep=1pt] at (59,-81) {\small $V_{20}$};
\node[inner sep=1pt] at (68,-77.5) {\small $V_{210}$};
\node[inner sep=1pt] at (92.5,-76.5) {\small $V_{21}$};
\node[inner sep=1pt] at (56.5,-91) {\small $V_2$};
\end{tikzpicture}
\end{document}
